\subsection{基域的改变}

设 \(K_{1}\) 是 K 的交换扩域。对于 K 上的李代数 g，它为半单的充要条件是 \(g_{(K_{1})}\) 为半单；因为 \(\beta_{1}\) 是 \(g_{(K_{1})}\) 的 Killing 形式，它由 g 的 Killing 形式 \(\beta\) 将基域从 K 扩张到 \(K_{1}\) 得到；因此，\(\beta_{1}\) 非退化当且仅当 \(\beta\) 非退化（《代数》，第九章，§ 1，第 4 号，命题 3 的推论）。

若 \(\mathfrak{g}_{(\mathbf{K}_{1})}\) 是单的，则根据上文 g 是半单的，且不可能是两个非零理想的乘积，因而 g 是单的。另一方面，即使 g 是单的，\(\mathfrak{g}_{(\mathbf{K}_{1})}\)（它是半单的）也可能不是单的（练习 17 和 26(b)）。

设 \(\mathfrak{g}\) 是李代数，\(\mathfrak{r}\) 是其根基。则 \(\mathfrak{r}_{(\mathbf{K}_1)}\) 是 \(\mathfrak{g}_{(\mathbf{K}_1)}\) 的根基（§ 5，第 6 号）。因此，若 \(\mathfrak{s}\) 表示 \(\mathfrak{g}\) 的幂零根，则 \(\mathfrak{g}_{(\mathbf{K}_1)}\) 的幂零根为 \([\mathfrak{g}_{(\mathbf{K}_1)},\mathfrak{r}_{(\mathbf{K}_1)}] = [\mathfrak{g},\mathfrak{r}]_{(\mathbf{K}_1)} = \mathfrak{s}_{(\mathbf{K}_1)}\)。由此可知，\(\mathfrak{g}\) 是约化的当且仅当 \(\mathfrak{g}_{(\mathbf{K}_1)}\) 是约化的。

设 \(\mathfrak{g}\) 是李代数，\(\mathfrak{h}\) 是其子代数。回忆，\(\mathfrak{h}\) 的一个表示是半单的，当且仅当由 \(\mathfrak{h}_{(\mathrm{K}_1)}\)（基域扩张到 \(\mathrm{K}_1\)）得到的表示是半单的。因此，\(\mathfrak{h}\) 在 \(\mathfrak{g}\) 中约化，当且仅当 \(\mathfrak{h}_{(\mathrm{K}_1)}\) 在 \(\mathfrak{g}_{(\mathrm{K}_1)}\) 中约化。

现在设 \(K_{0}\) 是 K 的子域，且 \([K:K_{0}]\) 有限。令 g 为李代数，\(g_{0}\) 为把基域从 K 限制到 \(K_{0}\) 后由 g 得到的（有限维）李代数。g 的每个交换理想都是 \(g_{0}\) 的交换理想；反之，若 \(a_{0}\) 是 \(g_{0}\) 的交换理想，则 g 中包含 \(a_{0}\) 的最小 K-向量子空间是 g 的交换理想；因此 g 是半单的当且仅当 \(g_{0}\) 是半单的。若 \(g_{0}\) 是单的，显然 g 是单的。反过来，设 g 是单的。我们证明 \(g_{0}\) 是单的。令 \(a_{0}\) 是 \(g_{0}\) 的一个单分量。对所有 \(\lambda \in K^{*}\)，\(\lambda a_{0}\) 是 \(g_{0}\) 的理想，并且

\[
\left[ a _ {0}, \lambda a _ {0} \right] = \lambda \left[ a _ {0}, a _ {0} \right] = \lambda a _ {0} \neq \{0 \},
\]

因而 \(\lambda a_0 \supset a_0\)，并由 \(\lambda a_0 = a_0\)（因为 \(\dim_{K_0}(\lambda a_0) = \dim_{K_0}a_0\)）可知该等式成立。现在，\(g\) 中由 \(a_0\) 生成的向量子空间是 g 的非零理想，因而等于 \(g\)。于是 \(g\) = \(g = a_0\)，这就证明了断言。

